% ===================================================================
% 00_nomenclature.tex -- Nomenclature (writer C1). ASCII only.
% ===================================================================
\chapter*{Nomenclature}
\addcontentsline{toc}{chapter}{Nomenclature}
\phantomsection\label{ch:nomenclature}

This list fixes the meaning of every symbol, abbreviation, provenance label and evidence class used in the report.
Values given here are those of the device and of the model as used throughout; the chapter or section in which a quantity is
derived is given in the last column where applicable. All currents are per micrometre of channel width
(A/$\mu$m); the simulated two-dimensional cross-section has a width of \SI{1}{\micro\metre}
(\cref{sec:inp-geometry}), and the total current of the \SI{290}{\micro\metre} wide device is 290 times the quoted value.

\section*{Device and bias symbols}
{\small
\begin{longtable}{@{}p{0.17\linewidth}p{0.53\linewidth}p{0.24\linewidth}@{}}
\toprule
Symbol & Meaning & Unit / value \\
\midrule
\endfirsthead
\toprule
Symbol & Meaning & Unit / value \\
\midrule
\endhead
\bottomrule
\endfoot
$\tIWO$, $t$ & IWO channel thickness (nominal film thickness of the measured series) & nm; 2.0, 6.3, 13.2, 31.8 \\
$L$ & channel length (gap between the Pd source and drain) & \SI{20}{\micro\metre} \\
$W$ & physical channel width & \SI{290}{\micro\metre} \\
$\Vg$, $\Vd$ & gate and drain voltage (source grounded) & V; $\Vd = \SI{0.7}{V}$ for all transfer curves \\
$\Id$, $\Ig$, $\Is$ & drain, gate and source terminal currents & A/$\mu$m \\
$\Cox$ & gate-stack capacitance per area (\HfO{} \SI{15}{nm} + \AlO{} \SI{2}{nm} in series) & \SI{8.955e-7}{\farad\per\centi\metre\squared} \\
$\EOT$ & equivalent SiO$_2$ thickness of the stack & \SI{3.86}{nm} \\
$q/\Cox$ & gate-voltage equivalent of a sheet charge of $10^{12}\cmsq$ & \SI{0.179}{V} \\
$T$ & lattice temperature (not recorded for the measurement; assumed) & \SI{300}{K} \\
$\kT$ & thermal energy at \SI{300}{K} & \SI{25.85}{meV} \\
\end{longtable}}

\section*{Extracted metrics}
The same extractor (\file{scripts/extract_metrics.py}) is applied to measured and simulated curves; the definitions are
derived in \cref{sec:der-metrics}.
{\small
\begin{longtable}{@{}p{0.17\linewidth}p{0.53\linewidth}p{0.24\linewidth}@{}}
\toprule
Symbol & Meaning & Unit \\
\midrule
\endfirsthead
\toprule
Symbol & Meaning & Unit \\
\midrule
\endhead
\bottomrule
\endfoot
$\Vthcc$ & constant-current threshold voltage, $\Vg$ at $\Id = 10^{-9}\Aum$ & V \\
$\Vthlin$ & linear-extrapolation threshold voltage (tangent at the transconductance maximum) & V \\
$\gm$, $\gmmax$ & transconductance $\partial\Id/\partial\Vg$ and its maximum & A/(V\,$\mu$m) \\
$\SSfix{I_1}{I_2}$ & fixed-current subthreshold swing: gate-voltage span per decade between the currents $I_1$ and $I_2$ (A/$\mu$m) & mV/dec \\
$\SScc$ & fixed-current swing between $10^{-10}$ and $10^{-8}\Aum$ & mV/dec \\
$\SSmin$ & minimum local swing of a curve; not a reliable metric (grid-phase and floor-window artefacts, \cref{sec:num-convergence}); reported for history only & mV/dec \\
$\muFE$ & linear-regime field-effect mobility, $L\,\gmmax/(W\Cox\Vd)$ & \cmVs \\
$\musat$ & saturation-formula mobility, $(2L/W\Cox)\,(\mathrm{d}\sqrt{\Id}/\mathrm{d}\Vg)^2_{\max}$, applied at $\Vd = \SI{0.7}{V}$ & \cmVs \\
$\Ion$ & drain current at $\Vg = \SI{3}{V}$ & A/$\mu$m \\
$\Ioff$, off-floor & median measured current in the off-band $-2 \le \Vg \le -0.5$\,V & A/$\mu$m \\
active region & measured points with $\Id > 5\times$ the off-floor (all points for $t \ge \SI{20}{nm}$); used for the fit residual & -- \\
active RMSE & root-mean-square of $\log_{10}\Id^{\mathrm{sim}} - \log_{10}\Id^{\mathrm{meas}}$ over the active region & dec \\
\end{longtable}}

\section*{Material and model parameters}
{\small
\begin{longtable}{@{}p{0.17\linewidth}p{0.53\linewidth}p{0.24\linewidth}@{}}
\toprule
Symbol & Meaning & Unit / section \\
\midrule
\endfirsthead
\toprule
Symbol & Meaning & Unit / section \\
\midrule
\endhead
\bottomrule
\endfoot
$\Eg$ & band gap (bulk value plus confinement increment) & eV, \cref{sec:par-eg} \\
$\chiIWO$ & electron affinity & eV, \cref{sec:par-chi} \\
$\dEc$, $\dEg$ & confinement shift of the conduction-band edge; equal to the gap increment ($\Ev$ unchanged) & eV, \cref{sec:par-dec} \\
$\mstar$ & electron effective mass (enters only $\Nc$) & $m_0$, \cref{sec:par-mstar} \\
$\Nc$, $\Nv$ & effective conduction- and valence-band densities of states & \cmcu, \cref{sec:par-nc,sec:par-nv} \\
$\epsIWO$ & relative permittivity of IWO & 9.3, \cref{sec:par-eps-iwo} \\
$\Nt$ & integrated acceptor-like tail density, $\NTA\WTA$ & \cmcu, \cref{sec:par-tail} \\
$\NTA$, $\WTA$ & tail DOS at $\Ec$ and its characteristic energy ($g_{\mathrm{TA}} = \NTA\exp[(E-\Ec)/\WTA]$) & \si{cm^{-3}.eV^{-1}}, eV \\
$\NGA$, $\EGA$, $\WGA$ & amplitude, depth below $\Ec$ and width of a Gaussian acceptor band & \si{cm^{-3}.eV^{-1}}, eV, eV \\
$\Dit$ & interface acceptor sheet at the IWO/\AlO{} interface (regularized in a \SI{0.25}{nm} layer) & \si{cm^{-2}.eV^{-1}}, \cref{sec:par-dit} \\
$\Qf$ & fixed front-interface charge density (in units of $q$) & \cmsq, \cref{sec:par-qf} \\
$\Qb$ & fixed charge on the exposed back (IWO/air) surface & \cmsq, \cref{sec:par-backq} \\
$\Ndeff$ & effective ionized background donor density (not the W concentration) & \cmcu, \cref{sec:par-nd} \\
$\muband$ & constant band mobility of a film (fitted per thickness) & \cmVs, \cref{sec:par-mu} \\
MUN & ATLAS constant electron mobility, $\muband(1-\Dsr/t)^2$ & \cmVs \\
$\Dsr$ & surface-roughness length of the roughness factor & \SI{0.287}{nm}, \cref{sec:par-rough} \\
$\Rsd$ & source/drain series resistance (zero in the model: ideal Ohmic contacts) & $\Omega\,\mu$m, \cref{sec:par-contacts} \\
$\Ea$ & activation energy of $\Id$ at fixed $\Vg$ & meV, \cref{sec:der-activation} \\
tmu, TMUN & ATLAS exponent of the constant-mobility temperature law $\mu(T) = \mu(300)(T/300)^{-\mathrm{TMUN}}$ & 1.5 (default, \cref{sec:par-temperature}) \\
NUMA/NUMD & number of energy levels used by ATLAS to integrate the continuous acceptor/donor DOS & 96/48 or 384/192, \cref{sec:der-dos-discretization} \\
\end{longtable}}

\section*{Abbreviations}
{\small
\begin{longtable}{@{}p{0.17\linewidth}p{0.77\linewidth}@{}}
\toprule
Abbreviation & Meaning \\
\midrule
\endfirsthead
\toprule
Abbreviation & Meaning \\
\midrule
\endhead
\bottomrule
\endfoot
IWO & tungsten-doped indium oxide (about 2 at.\,\% W in \InO{}; composition definition not determined from the available data) \\
TFT & thin-film transistor \\
TCAD & technology computer-aided design (here: Silvaco ATLAS 5.28.1.R drift-diffusion) \\
ALD & atomic layer deposition \\
RTA & rapid thermal annealing \\
DOS & density of states \\
DFT, PBE & density functional theory; Perdew-Burke-Ernzerhof functional \\
EMA & effective-mass approximation \\
SP & Schr\"odinger-Poisson (1-D surrogate of specialist S3) \\
MTR & multiple trapping and release \\
SRH & Shockley-Read-Hall recombination \\
KCL & Kirchhoff current law check, $\max|\Id+\Is+\Ig|$ \\
CNL & charge-neutrality level \\
TLM & transfer-length method (not available) \\
PBS & positive bias stress \\
SI & supporting information of \citet{Januar2026} (not bundled) \\
V0, V1 & earlier per-thickness package (\file{IWO_ATLAS_Model_claude}, 34 launches) and the present physics-constrained package \\
Astra & earlier independent package (\file{astra_IWO_PRIORITY_THREE_20260912}), used for cross-checks \\
S1--S6 & specialist analyses of 2026-09-25 (S1 electrostatics, S2 transport, S3 confinement, S4 defects, S5 contacts/validation, S6 numerics/predictions) \\
\end{longtable}}

\section*{Run identifiers}
Every simulated number carries its run identifier \run{NNNN} ($\mathrm{NNNN} = 0001\ldots0052$), the name of the folder
\file{results/runs/run_NNNN_<film>_<label>} in the package; the film key is written \texttt{2p0}, \texttt{6p3},
\texttt{13p2}, \texttt{31p8}. The complete list is \cref{tab:run-catalogue} (\cref{ch:app-runs}). Run identifiers of the
earlier V0 package are always written ``V0 run\_NNNN'' to avoid confusion. Campaign codes of 2026-09-25 map onto run
identifiers as follows: A1 = \run{0031}, A2 = \run{0032}, A3 = \run{0030}, A4 = \run{0033}, A5 = \run{0035},
A6 = \run{0034}, A7 = \run{0036}; B0 = \run{0037}, B1--B3 = \run{0038}--\run{0040}, B4--B7 = \run{0048}--\run{0051}
(isolation tests), B8--B10 = \run{0041}--\run{0043} ($\Id$--$\Vd$ predictions), B11--B14 = \run{0044}--\run{0047}
(elevated temperature); C1 = \run{0052}.

\section*{Provenance labels}
Every parameter carries one label; the labels are those of \file{config/iwo_material_model.yaml} and the master table
\cref{tab:provenance-master}.
{\small
\begin{longtable}{@{}p{0.30\linewidth}p{0.64\linewidth}@{}}
\toprule
Label & Meaning \\
\midrule
\endfirsthead
\toprule
Label & Meaning \\
\midrule
\endhead
\bottomrule
\endfoot
\prov{MEASURED\_TARGET\_DEVICE} & measured on, or stated for, the target devices of \citet{Januar2026} or the user's schematic and workbook \\
\prov{LITERATURE\_IWO} & value reported for IWO, but for a different device or process \\
\prov{MEASURED\_RELATED\_MATERIAL} & measured on a related material (e.g.\ bulk \InO{}) \\
\prov{DFT\_PURE\_IN2O3\_PROXY} & computed by DFT for pure \InO{} slabs and used as a proxy for IWO \\
\prov{LITERATURE\_IN2O3\_PROXY} & literature value for an \InO{} device used as a proxy (e.g.\ $\Dit$ of \HfO{}/\InO{}) \\
\prov{FITTED} & adjusted to the measured transfer curves of this study (calibration) \\
\prov{ASSUMED} & chosen without direct evidence; the choice and its sensitivity are stated \\
\prov{CALCULATED} & computed from other labelled quantities by a stated formula \\
\end{longtable}}

\section*{Evidence classes and verdicts}
{\small
\begin{longtable}{@{}p{0.17\linewidth}p{0.77\linewidth}@{}}
\toprule
Class & Meaning \\
\midrule
\endfirsthead
\toprule
Class & Meaning \\
\midrule
\endhead
\bottomrule
\endfoot
\evDATA & a measured curve or a number extracted from it \\
\evNUM & converged ATLAS result; a statement about the model, not about the device \\
\evCAL & agreement with the same curve that was used to fit the parameters (calibration, not validation) \\
\evOOS & out-of-sample test designed before the outcome was known \\
\evPOST & test designed or evaluated after the outcome was known \\
\evPRED & registered prediction, not yet tested against measurement \\
\evSURR & result of a specialist surrogate model (1-D Poisson, SP, charge sheet), not of ATLAS \\
\evTHEORY & analytic argument \\
\evLIT & statement of the literature, with page \\
\evCORR & correction of an earlier claim of this project \\
\end{longtable}}
Verdicts are \verdict{demonstrated}, \verdict{strongly supported}, \verdict{plausible, unverified},
\verdict{rejected} and \verdict{not determined}. ``Calibrated'' (\evCAL) and ``validated'' (\evOOS{} passed) are never used
interchangeably.
